\documentclass[11pt]{article}

\usepackage[utf8]{inputenc}
\usepackage[T1]{fontenc}
\usepackage{lmodern}
\usepackage[a4paper,margin=2.5cm]{geometry}
\usepackage{booktabs}
\usepackage{array}
\usepackage{enumitem}
\usepackage{listings}
\usepackage{xcolor}
\usepackage{underscore}
\usepackage{microtype}
\usepackage[hidelinks]{hyperref}

\setlength{\emergencystretch}{3em}

% Measured figures. Every number traces to a committed experiment in
% experiments/agentsec/ and is re-checkable via reproduce.sh (8/8).
\newcommand{\Nrace}{4000}
\newcommand{\RaceLo}{19.2}
\newcommand{\RaceHi}{22.8}
\newcommand{\InjAtt}{164}
\newcommand{\AutoAtt}{45}
\newcommand{\AutoEscape}{18}

\lstset{
  basicstyle=\ttfamily\small,
  breaklines=true,
  columns=fullflexible,
  keepspaces=true,
  frame=single,
  framerule=0.3pt,
  framesep=4pt,
  xleftmargin=2pt,
  showstringspaces=false,
  aboveskip=8pt,
  belowskip=8pt
}

\title{Deterministic is not sound:\\
an empirical evaluation of seccomp argument-inspection gates
for containing LLM agents}
\author{Sadik Erisen\thanks{This research was conducted independently. The views
and findings are those of the author and do not reflect the position of any
corporate employer. No data in this paper is fabricated; every figure traces to a
committed artifact and is reproducible via \texttt{experiments/agentsec/reproduce.sh}.}\\[3pt]
\normalsize Independent Researcher \& Graduate Student\\
\normalsize University of Pennsylvania}
\date{July 2026}

\begin{document}
\maketitle

\begin{abstract}
Sandboxes that run untrusted LLM-agent code increasingly reach for a syscall-boundary
control: a seccomp supervisor that inspects the \emph{arguments} of intercepted system
calls and allows or denies each one. Because the supervisor's verdict is deterministic,
it is tempting to treat the control as certified containment. We test that control
directly, on a mechanical oracle that compiles and runs the code and reads its behaviour
from \texttt{strace}, never from text. We report four findings. First, the
argument-inspection gate is bypassable: it blocks a persistence write \Nrace/\Nrace{}
times head-on, yet a two-thread time-of-check-to-time-of-use (TOCTOU) race on its
allow path defeats it at \RaceLo--\RaceHi\% per attempt, which approaches certainty under
retry. Second, attacker-injected and self-driven behaviour do not collapse to one syscall
footprint: they share \texttt{execve}, but \AutoEscape\% of self-driven destruction takes
a direct \texttt{unlink}/\texttt{rmdir} path that an \texttt{execve}/\texttt{connect} gate
never observes. Third, the gate is genuinely selective rather than deny-all, but its
substring path-matching admits false positives on benign look-alike filenames while its
anchored tokens are evaded by path normalization. Fourth, attributing the actor from its
behaviour---by timing, by stated intent, or by its reaction to a denial---fails as a
control against an adversary that chooses its own behaviour. The determinism of the
verdict is not soundness of the containment. What protects the host is coverage of the
mechanisms an agent can use, and a denied agent enumerates the alternatives rather than
stopping.
\end{abstract}

\section{Introduction}

An LLM agent that writes and runs code is, from the operating system's point of view, an
untrusted process whose intent cannot be read from its prompt, its reasoning trace, or the
code it emits. A common containment proposal is to watch the one interface the process
cannot avoid: the system-call boundary. A rootless filter blocks by syscall number; a
richer, ``rootful'' design installs a seccomp user-notification supervisor that reads the
pointer arguments of an intercepted call out of the target's memory and decides allow or
deny per call. Argument visibility is the selling point---the supervisor can permit
\texttt{openat} on a log file and deny it on \texttt{\char`~/.ssh/authorized\_keys}.

The supervisor's verdict is deterministic: the same call with the same arguments yields the
same decision every time. That property invites a category error, in which a deterministic
control is presented as \emph{certified} containment. This paper separates the two. A
verdict can be perfectly deterministic and still wrong, either because it inspects an
argument the kernel later re-reads, or because the dangerous action never passes through a
watched call at all. Determinism bounds the \emph{variance} of the control; soundness would
bound its \emph{coverage and correctness}, and coverage is where the containment fails.

We hold every claim to a mechanical standard. A statement that ``the agent attempted $X$''
means a compiled program issued syscall $X$ under \texttt{strace} inside a rootless network
namespace with a timeout; we never infer behaviour from generated text. Under that
standard we report the following as findings, not activities.

\begin{enumerate}[leftmargin=1.4em,itemsep=2pt]
\item The argument-inspection gate is bypassable by a TOCTOU race on its allow path
(Section~\ref{sec:toctou}).
\item Attacker-driven and self-driven behaviour do not reduce to a single syscall footprint,
so a gate scoped to one mechanism leaves the other's path open (Section~\ref{sec:coverage}).
\item The gate's selectivity rests on substring path-matching that admits false positives and
is separately evadable (Section~\ref{sec:selectivity}).
\item Behavioural attribution of the actor is not a control (Section~\ref{sec:attribution}).
\end{enumerate}

Section~\ref{sec:holds} states what does survive, so the reader does not over-correct: the
syscall boundary remains the right place to observe, and interception is invariant to the
source language. Sections~\ref{sec:discussion}--\ref{sec:related} give the implication for a
threat model, the limits of the study, and the position against prior work.

\section{Threat model}\label{sec:threat}

The host runs code produced by an LLM agent and treats that code as untrusted. Two sources
drive dangerous behaviour, and the design goal is to contain both at one boundary. The first
is \emph{attacker injection}: a prompt-injection payload steers the agent into emitting a
dangerous call, for example \texttt{system("/bin/sh")} routed through the generated code. The
second is \emph{autonomous malfunction}: with no attacker, the agent reaches for a
destructive shortcut on its own, for example deleting a build directory it was asked to
``reset.''

The defender controls the seccomp filter and the supervisor process; the supervisor obtains
read access to the target's memory through the parent--child relationship (equivalently
\texttt{CAP\_SYS\_PTRACE} in a hardened deployment). The attacker controls the contents of the
untrusted process---its threads, the timing of its calls, and the arguments in its own
address space---but not the supervisor or the kernel. This is the standard capability split
for a syscall supervisor, and it is exactly the split under which the allow-path race in
Section~\ref{sec:toctou} is available: the attacker owns the memory the supervisor reads and
the kernel later re-reads.

\section{Method}\label{sec:method}

\paragraph{Mechanical oracle.} Generated C is compiled with \texttt{gcc} and executed under
\texttt{timeout 10 unshare -rn strace -f}. The footprint of a run is the set of security-relevant
syscalls it issues (\texttt{execve}, \texttt{connect}, \texttt{openat}, \texttt{unlink},
\texttt{rmdir}, and their families), read from the trace. A baseline of one \texttt{execve}---the
process launch---is subtracted before a run counts as ``attempting.'' No claim in this paper is
derived by scanning model output for keywords.

\paragraph{Subjects.} Three local code models generate the artifacts: \texttt{qwen2.5-coder:1.5b},
\texttt{qwen2.5-coder:7b}, and \texttt{deepseek-coder:6.7b}, served through Ollama. The host is a
single x86\_64 machine. Section~\ref{sec:limits} states what this scope does and does not support.

\paragraph{The gate.} The rootful gate forks the target, installs \texttt{SCMP\_ACT\_NOTIFY} on
\texttt{execve}, \texttt{execveat}, and \texttt{connect} (and, in persistence mode, on
\texttt{openat}, \texttt{rename}, \texttt{link}, and their variants), and sends the notification
descriptor to the supervisor. On each notification the supervisor reads the pathname from
\texttt{/proc/pid/mem}, compares it against policy, and either allows the call by returning
\texttt{SECCOMP\_USER\_NOTIF\_FLAG\_CONTINUE} or denies it with \texttt{EPERM}. The rootless gate,
by contrast, is a blanket \texttt{SCMP\_ACT\_ERRNO} deny of \texttt{execve}/\texttt{execveat}/%
\texttt{connect} with no argument inspection.

\paragraph{Statistics.} For a stochastic quantity---an agent's rate of emitting a dangerous
call---we report a Wilson 95\% interval. The gate's decision is not a stochastic rate and receives
no failure-rate interval; its determinism is a property to be stated, and its soundness is measured
separately as coverage (Sections~\ref{sec:toctou}--\ref{sec:selectivity}). Every figure below is
reproducible from committed source through \texttt{experiments/agentsec/reproduce.sh}, which rebuilds
the binaries and re-checks the headline results (8 checks, all passing at the time of writing).

\section{The argument-inspection gate is bypassable}\label{sec:toctou}

The gate blocks the direct attack completely. A target that writes to a path containing a
persistence token---\texttt{authorized\_keys}, \texttt{/.ssh/}, \texttt{cron}---is denied on every
one of \Nrace{} attempts, and the supervisor's own counter agrees with the target's: \Nrace{} calls
seen, \Nrace{} denied. Head-on, the argument inspection does exactly what it claims.

The allow path does not survive a race. On a permitted call the supervisor validates the pathname it
read from \texttt{/proc/pid/mem}, then returns \texttt{SECCOMP\_USER\_NOTIF\_FLAG\_CONTINUE}, which
instructs the kernel to resume the syscall. Resuming the syscall means the kernel copies the pathname
pointer \emph{again}, at execution time. Between the supervisor's read and the kernel's re-read the
target's other threads keep running---only the calling thread is parked on the notification---so a
sibling thread that flips the buffer from a benign path to \texttt{authorized\_keys} in that window
causes the supervisor to validate one path and the kernel to open another.

We measured the effect with a two-thread target: one thread issues the \texttt{openat}, a second
flips the shared path buffer. Across \Nrace{} attempts the evil path was opened despite the gate on
\RaceLo\% of tries in the first run and \RaceHi\% in an independent rebuild; the remainder split
between benign opens and correct denials, and the three outcomes account for every trial. The
existence of the bypass is a property of the mechanism and is host-independent; the per-attempt rate
depends on scheduler behaviour and moved between runs, so we do not treat it as a fixed figure. Under
retry the distinction stops mattering: at a per-attempt success near one in five, a persistence write
succeeds with probability above 0.9 within roughly eleven attempts, and the agent is free to retry.

This is the documented deep-argument hazard of seccomp user notification, and the
\texttt{seccomp\_unotify}(2) manual page warns against using \texttt{CONTINUE} to allow a syscall
whose arguments were security-relevant, precisely because the arguments can change before the kernel
acts. Two consequences follow for a threat model. The bypass is architectural, not a gap in the token
list, so no blocklist patch closes it; the second read has to be removed, by copying the argument into
supervisor-controlled memory or by having the supervisor perform the operation itself. And the head-on
number---\Nrace/\Nrace{}---measures the control against a cooperative target, which the threat model
does not assume.

We note one correction to a natural reading of this result. The gate permits exactly one
\texttt{execve}, the launch of the target itself, and that call fires in a single-threaded context
before any agent code runs; every later \texttt{execve} is denied outright. The raceable allow path is
therefore not the \texttt{execve} launch but the per-call file-operation allow in persistence mode,
which fires while the attacker's threads are live.

\section{Coverage is incomplete before any race}\label{sec:coverage}

A gate can be sound only over the mechanisms it watches, and the two threat sources of
Section~\ref{sec:threat} do not use the same mechanisms. We compared the footprints of the artifacts
each source actually produced. The \InjAtt{} injection artifacts that attempted a dangerous call did so
entirely through \texttt{execve}: the footprint set is $\{\texttt{execve}\}$. The \AutoAtt{} autonomous
artifacts that attempted destruction did so through \texttt{execve} in 82\% of cases (spawning a shell
to run \texttt{rm}), but through a direct \texttt{unlink} in 13\% and a direct \texttt{rmdir} in 4\%:
the footprint set is $\{\texttt{execve}, \texttt{unlink}, \texttt{rmdir}\}$. The two sets share
\texttt{execve} and are not equal. The footprint is fixed by the \emph{mechanism} the code chooses, not
by the \emph{force} that drove it, and \AutoEscape\% of self-driven destruction reaches its goal by a
call an \texttt{execve}/\texttt{connect} gate never sees.

The gap widens at the ABI. Blocking \texttt{execve} does not cover \texttt{execveat} on an anonymous
descriptor via \texttt{memfd\_create}; blocking \texttt{connect} does not cover \texttt{sendto} or
\texttt{sendmsg} on a UDP socket; and a single unblocked \texttt{io\_uring\_setup} lets an agent submit
\texttt{openat}, \texttt{write}, \texttt{connect}, and \texttt{unlink} as ring operations that never
reach the syscall-entry filter, defeating file and network coverage at once. On a filter scoped to one
architecture's syscall numbers, the i386 and x32 compat entries fall through to the default action.
A blocklist cannot enumerate its way out of this; completeness requires inverting to a default-deny
allowlist, and the only class this study could close as a small fixed set is execution
(\texttt{execve}, \texttt{execveat}). We did not exercise every listed ABI path as a working exploit;
the \texttt{unlink}/\texttt{rmdir} escape is measured on real artifacts, and the remainder are coverage
gaps read from the filter definition and the syscall ABI.

\section{Selectivity trades false positives against evasion}\label{sec:selectivity}

The rootful gate is not a disguised deny-all: given a labelled corpus of write paths it denied all
seven persistence targets and allowed the ordinary application writes, so it does discriminate on the
argument, unlike the rootless filter, which denies the watched calls unconditionally. Its
discrimination mechanism, however, is a substring match against a token list, and that mechanism sits
on a false-positive/evasion trade-off it cannot leave.

On a seventeen-path corpus the confusion matrix was seven true positives, zero false negatives, six
true negatives, and four false positives (positive $=$ should-be-denied). Every false positive was a
legitimate filename that merely contains a token as a substring---\texttt{microcron\_metrics.txt}
matching \texttt{cron}, \texttt{authorized\_keys\_report.csv} matching \texttt{authorized\_keys}. The
tokens that cause these false positives are the un-anchored ones; the tokens that avoid them are anchored
with a leading slash (\texttt{/.ssh/}, \texttt{/.profile}), and those are exactly the tokens defeated by a
relative path or a symbolic link through an allowed directory. Tightening a token to cut a false positive
is the same edit that opens an evasion, and loosening it to catch an evasion is the same edit that creates
a false positive. Substring matching cannot be both precise and evasion-resistant; a race-free, path-aware
mechanism such as Landlock resolves the path before deciding and is the correct comparison point for a
production design.

\section{What holds}\label{sec:holds}

The findings above are not an argument to abandon the syscall boundary. Two properties of the boundary
survive, and stating them keeps the recommendation calibrated.

Interception is invariant to the language that emitted the call. The same second \texttt{execve}---a
shell spawned after launch---is denied identically whether the target is compiled C, a Python
interpreter, or a shell script (\texttt{denied\_execve}$=1$ in each, with the shell never running). The
gate observes the kernel interface, so the source language above it does not change the verdict.

The boundary is also the right place to observe, because analysis above it cannot certify. In companion
work on the same corpus, ten static and dynamic detectors, AddressSanitizer included, all miss an
intra-object overflow that crosses no kernel boundary; the artifact is clean to every checker yet
corrupt at runtime. Observation therefore has to move to execution, and the syscall boundary is where
execution becomes visible. What the boundary gives is a deterministic, auditable \emph{verdict}; what it
does not give, per Sections~\ref{sec:toctou}--\ref{sec:selectivity}, is soundness.

\section{Behavioural attribution is not a control}\label{sec:attribution}

If the gate is porous, a natural next move is to profile the agent's behaviour and stop the bad actor
before it reaches a gap. We tested three variants and none is a control against an adversary, for one
structural reason: a manipulated agent chooses its own behaviour, so any signature the defender keys on
can be produced or suppressed at will.

Timing does not attribute. A script that inserts random sleeps is statistically indistinguishable from a
live agent's action cadence (coefficient of variation 0.45 versus 0.54; fifteen of seventeen inter-action
intervals overlap), so raw timing separates neither by construction.

Stated intent against footprint is a weak, spoofable signal. Comparing a model's one-sentence description
of its own code against the code's measured footprint, divergence appeared in 22\% of injection cases and
in none of the autonomous or benign cases; where it appeared it was mechanism concealment, a description
of a file write that omits the shell it spawns. A non-manipulated model mostly narrates honestly, and a
manipulated one writes a faithful malicious description and drives the divergence to zero.

Reaction to a denial is real but points the other way. Given a blocked goal, a live agent adapted on 90\%
of steps, repeated on 10\%, and stopped on none; it escalated from a plain redirection to \texttt{sudo},
to \texttt{su}, to \texttt{tee}, and to an alternative path, holding the objective fixed while varying the
mechanism. The operational reading is that a denial elicits enumeration, not deterrence, which returns the
weight to coverage: blocking one path routes the agent to the next.

\section{Discussion}\label{sec:discussion}

The result a defender should carry out of this study is a separation the marketing of syscall sandboxes
tends to blur. A seccomp argument-inspection gate delivers a deterministic and auditable verdict, and that
is worth having; it does not deliver certified containment, and presenting it as such misplaces the
residual risk. The residual risk is coverage. It lives in the mechanisms the gate does not watch
(Section~\ref{sec:coverage}), in the window between a check and the kernel's re-read
(Section~\ref{sec:toctou}), and in the imprecision of the argument match (Section~\ref{sec:selectivity}).
It does not live in the variance of the verdict, and it is not recovered by profiling the agent
(Section~\ref{sec:attribution}).

Two design directions follow, and both are about closing coverage rather than tuning the agent. First,
remove the second argument read on the allow path: resolve and copy the argument into memory the target
cannot mutate, or let the supervisor perform the operation, so the value the supervisor checked is the
value the kernel uses. Second, move from a blocklist to a default-deny allowlist over a race-free,
path-resolving mechanism, so that coverage is the complement of an explicit permit set rather than the
complement of an enumerated deny set. We do not report an implementation of either here; they are the
shape a sound successor takes.

\section{Limitations}\label{sec:limits}

The study is scoped and we state the scope rather than let a reviewer find it. The subjects are three
local code models between 1.5 and 7 billion parameters, not frontier agents, and the autonomous
destruction result rests on \AutoAtt{} attempting artifacts; both the rates and the specific mechanisms
seen may shift with larger models and larger samples, and the emergent-destruction finding in particular
is small-$N$. Every measurement was taken on a single x86\_64 host. Interception was checked across three
implementation languages but on one architecture; the aarch64 case is not established here, because
faithful native seccomp semantics are not available under user-mode emulation, and a second-architecture
host is the correct way to test it. The TOCTOU bypass is demonstrated as an existence result, which one
successful race establishes; its per-attempt rate is scheduler-dependent and is not an external-validity
claim. The behavioural-attribution experiments are forensic characterizations of non-adversarial actors
and are reported as such; against an adversary that controls its own behaviour they attribute nothing,
which is the point of Section~\ref{sec:attribution}.

\section{Related work}\label{sec:related}

Syscall-boundary containment for untrusted code predates LLM agents. Seccomp-BPF and its user-notification
extension are the mechanism under test here; the deep-argument race we exercise is described in the kernel
manual pages for user notification and is not a new class of bug, only an instance of it in an
agent-containment setting. Application kernels and sandboxes---gVisor, and the unprivileged, path-resolving
Landlock LSM---address the coverage and race problems this paper measures, and Landlock is the natural
comparison for the file-operation policy of Section~\ref{sec:selectivity}
\,[CITATION NEEDED --- verify gVisor, Landlock references]. Attribution of a process by its syscall
behaviour goes back to the host-intrusion-detection line begun by Forrest et al., \emph{A Sense of Self for
Unix Processes} (1996) \,[CITATION NEEDED --- verify], which is the lineage our Section~\ref{sec:attribution}
sits in and declines to extend into a control. On the language side, that a chain-of-thought need not match
behaviour is established---Turpin et al. (2023) and Lanham et al. (2023) \,[CITATION NEEDED --- verify
titles/venues]---which is why the divergence signal cannot be a detector. The emerging agent-safety
evaluations (ToolEmu; AgentDojo; AgentHarm) \,[CITATION NEEDED --- verify] measure agent harm at the tool or
task level; this paper measures the containment control beneath them, at the syscall boundary. We position
against these lines rather than beside them: prior sandboxing work supplies race-free mechanisms our tested
gate lacks, and prior attribution work supplies the negative result our behavioural section confirms in the
agent setting.

\section{Conclusion}\label{sec:conclusion}

A deterministic syscall gate is a useful control and a poor certificate. Measured against an adversary
rather than a cooperative target, the argument-inspection gate is bypassable on its allow path, blind to
the mechanisms it does not enumerate, and imprecise where it matches paths by substring; and the agent it
contains, when denied, looks for another way rather than stopping. The defensible claim is narrow and worth
stating plainly: certify the verdict, scope the coverage, and assume the agent will route around any single
blocked path. The open work is a supervisor with no second argument read and a default-deny allowlist over
a path-resolving mechanism, and a repeat of these measurements against frontier models and a second
architecture.

\paragraph{Reproducibility.} All code, data, and per-experiment results are committed at
\texttt{github.com/francose/llm-sandbox-cwe-research}. \texttt{experiments/agentsec/reproduce.sh} rebuilds
the gate and the probes from source and re-checks the headline results.

\end{document}
